% Operations letter to terminal management (~1.5 pages).
% Compile: pdflatex airport_authority_letter.tex
% Style: plain management English; short sentences; real KPIs only.
\documentclass[11pt]{article}
\usepackage[margin=1in]{geometry}
\usepackage{setspace}
\usepackage[hidelinks]{hyperref}
\pagestyle{plain}
\onehalfspacing

\begin{document}

\begin{center}
{\Large\bfseries Optimizing Terminal Passenger Charging Services:\\
A Data-Driven Strategic Blueprint}\\[0.6em]
{\normalsize Memorandum to the Airport Terminal Management Committee}\\[0.25em]
{\small Team GoHiMCM \quad HiMCM 2019 Problem A study pack}\\[0.15em]
{\small Evidence files: \texttt{optimal\_kiosk\_sizing.csv},
\texttt{thunderstorm\_stress\_comparison.csv}}
\end{center}

\vspace{0.8em}
\noindent
Dear Members of the Committee,

\medskip
\noindent
Passengers treat a charged phone as part of the journey product.
Dead batteries raise missed-connection risk and desk complaints.
We ran a local discrete-event model of charging kiosks for one
$16$-hour day.
We report only measured simulation outputs.

\section*{What the model shows}
CapEx is capped at $\$80{,}000$.
Floor area is capped at $40\,\mathrm{m}^2$.
A rush-day search evaluated $216$ packs.
Seven packs remain on the Pareto front.
The recommended knee pack is
\textbf{8 fast kiosks / 16 slow lockers / 0 power-bank banks}.
CapEx is $\$35{,}200$.
That CapEx uses $44\%$ of the CapEx ceiling.
Floor use is $38.4\,\mathrm{m}^2$.
Rush mean wait is $0.0115$ minutes.
Rush reneging probability is $0.00$.

We injected a thunderstorm delay at $17{:}00$.
Forty percent of later departures gained $2$--$4$ hours of delay.
Phone drain doubled for delayed travelers.
The run has $n=621$ arrivals.
The empirical pack is $12/15/2$ at CapEx $\$49{,}000$.
The knee cut mean wait from $4.591$ to $2.347$ minutes.
That is a $48.9\%$ wait cut.
Reneging probability fell from $0.0229$ to $0.0114$.
That is a $50.3\%$ reneging cut.
CapEx fell by $28.2\%$.
Peak wait-queue length stayed at $6$ for both packs.

\section*{Standard operating procedures (SOP)}
\begin{enumerate}
  \item \textbf{Three-minute walk grid.}
  Place charging kiosks so every gate cluster is within a three-minute
  walk.
  Prefer corridor nodes that serve two gates.
  Do not stack all fast kiosks at one retail island.

  \item \textbf{Hardware mix for Phase 1.}
  Procure pack $8/16/0$ first.
  Keep $\$44{,}800$ of CapEx headroom under the $\$80{,}000$ cap.
  Use headroom only after a second storm-seed review.

  \item \textbf{Free first 15 minutes on the fast kiosk.}
  Post a clear tariff: first $15$ minutes free on the fast kiosk;
  meter thereafter.
  The free window raises turnover and cuts abandoned waits.
  Keep slow lockers as the long-dwell product.

  \item \textbf{Storm playbook after $17{:}00$.}
  When Ops declares a weather delay, open a temporary phone-charging
  lane at the largest hold room.
  Staff one attendant to load-balance fast kiosks and slow lockers.
  Publish queue status on gate screens every $10$ minutes.

  \item \textbf{Power-bank policy.}
  Phase 1 does not require power-bank banks at the knee pack.
  If retail demands banks, add them only where gate dwell exceeds
  $90$ minutes and CapEx remains under the cap.

  \item \textbf{Monthly KPI board.}
  Track mean wait, reneging probability, and peak queue length.
  Trigger a re-size review if reneging probability exceeds $5\%$
  for two consecutive peak weeks.
\end{enumerate}

\section*{Ask of the Committee}
Approve Phase~1 procurement of pack $8/16/0$.
Authorize the free-$15$-minute fast-kiosk tariff pilot for $90$ days.
Direct Ops to draft the $17{:}00$ storm lane checklist before the next
severe-weather season.

\medskip
\noindent
We can supply the locked CSV tables and figures on request.
All figures are $300$\,DPI outputs from the same local Python runs.

\bigskip
\noindent
Respectfully,\\[0.6em]
GoHiMCM Study Team\\
Airport Charging Resilience Workstream

\end{document}
